\section{拟映射墙交叉（Yang Zhou）}%
\label{sec:Quasimap wall-crossing (Yang Zhou)}

本讲座中的主墙交叉定理与主空间构造取自~\cite{Zhou2022}。更早的亏格零与高亏格结果见~\cite{CFK2014,CFK2020}，稳定拟映射的基础见~\cite{CFKM2014}。

设 \(X\) 为五次三维流形。我们的目标是对 \(X\) 上的亏格零曲线计数，这个问题已被 Givental 与 Lian-Liu-Yau 解决。通过闭浸入 \(X \subset \P^4\)，我们可以写成
\[ \bar{M}_0(X, d) \subset \bar{M}_0(\P^4, d), \]
它是光滑的，并带有环面 \((\C^{\times})^5\) 的作用。若一个映射被该作用固定，则其定义域必有一些分支映到不动点，另一些分支构成 \(1\) 维轨道的 \(n :1\) 覆叠。特别地，不动点轨迹由一族图来标记。

在这一设定下，我们有量子 Lefschetz 原理，即
\[ [\bar{M}_0(X, d)]^{\on{vir}} = e(E) \cap [\bar{M}_0(\P^4, d)]. \]
其中，向量丛 \(E\) 定义为
\[ E = \pi_* f^* \mc{O}(5), \]
其中 \(\pi\) 是万有曲线，\(f\) 是万有稳定映射。对下述结果我们将给出一种不同的视角：若
\begin{align*}
    I(q, z) &= 1 + \sum_{d > 0} \frac{\prod_{k=1}^{5d} (5H+kz)}{\prod_{k=1}^d (H+kz)^5} q^d \in H^*(X)[z^{-1}] \ps{q} \\
    &= I_0(q) \cdot 1 + I_1(q) \frac{H}{z} + I_2(q) \frac{H^2}{z^2} + I_3(q) \frac{H^3}{z^3},
\end{align*}
则
\[ J (Q, z) = 1+ \sum_{d>0} \phi^{\alpha} \ab<\frac{\phi_{\alpha}}{z(z-\psi)}>_{0,1}^X Q^d = \frac{I}{I_0} \]
在镜面映射 \(Q = q \exp\ab(\frac{I_1}{I_0})\) 的意义下成立。 

\subsection{朴素计数}%
\label{sub:Naive counting}

我们将考虑从固定的 \(\P^1\) 出发的``映射'' \(\P^1 \to X \subset \P^4\)。为简单起见，我们假设 \(X = \ab(\sum x_i^5 = 0)\) 是 Fermat 五次流形。它由一组 \([f_0, \ldots, f_4]\) 给出，其中 \(f_i \in \C[u,v]\) 满足 \(\deg f_i = d\) 且 \(\sum f_i^5 = 0\)。要成为真正的态射，需要这些 \(f_i\) 没有公共零点，但这样得到的模空间不是本征的，因此我们放宽这一条件。于是我们考虑模空间
\[ QG(X, d) =  \ab\{ [f_0, \ldots, f_4] \middle\vert \Vectorstack{{ f_i \in \C[u,v] } { \deg f_i = d }}, \sum f_i^5 = 0, (f_0, \ldots, f_4) \neq (0,\ldots, 0) \} \bigg/ \C^{\times}. \]
它包含在相应的到射影空间的有理映射的模空间
\begin{align*}
    QG(\P^4, d) &\coloneqq  \ab\{ [f_0, \ldots, f_4] \mid f_i \in \C[u,v], \deg f_i = d, (f_0, \ldots, f_4) \neq (0, \ldots, 0) \} \\
    &\cong \P^{5d+5-1} 
\end{align*}
之中。这里有两个问题。第一，定义域被参数化了，因此我们需要考虑 \(\P^1\) 的自同构。第二，要得到稳定映射，我们需要对定义域作爆破。

首先，令 \(\C^{\times}\) 通过缩放 \(u\) 作用于 \(\P^1\)。考虑特殊的不动点轨迹 \(F_{\star} \subset QG(X, d)\)，其定义为
\[ F_{\star,d} = \ab\{ [f_0, \ldots, f_4] \text{ fixed and } (v=0) \text{ is not a common zero} \}. \]
因此，\(F_{\star,d} = \ab\{ [a_0 u^d, \ldots, a_4 u^d] \mid [a_0, \ldots, a_4] \in X \} \cong X\)。从这个视角，我们可以把 I-函数改写为
\begin{align*}
    I(q,z) &= \sum_{d > 0} \ab(\frac{[F_{\star, d}]^{\vir}}{e_{\C^{\times}}(N^{\vir})})^{\mr{PD}} q^d \\
    &= 1 + \sum_{d > 0} \ab( \frac{[X]}{e_{\C^{\times}}(N^{\vir})})^{\mr{PD}} q^d
\end{align*}
这由形变理论中的标准论证得到。它活在 \(H^*_{\C^{\times}}(X) \ps{q}[z^{-1}]\) 中。

\subsection{不严格的墙交叉}%
\label{sub:Wall-crossing without rigor}


其次，考虑一条结点曲线 \(C\)，它带有一个标记点外加一条带有 \(\C^{\times}\) 作用的尾，并设我们有映射 \(f \colon C \to X\)。若取极限
\[ \lim_{t \to 0} [f_0(t^{-1}u, v), \ldots, f_4(t^{-1}u, v)] = [a_0 u^{d_0}, \ldots, a_4 u^{d_0}], \]
我们可以看到泡泡的极限落在 \(F_{\star, d_0}\) 中，其中公共零点远离结点。于是我们可以用一个新标记点替换这条尾，它携带着 I-函数的某些项。对五次流形而言，这恰好是 \(I_1 H - I_0 \psi = [zI-z]_{+, z = -\psi}\)。在这一过程结束时，我们只剩下 \(C = \P^1\) 以及 I-函数的种种贡献。

注意，若考虑 \(\C^{\times}\) 作用于 \(\P^4\)，则 \(f \colon C \to \P^4\) 被固定当且仅当对所有 \(t \in \C^{\times}\)，存在 \(\varphi_t\) 使得
\begin{equation*}
\begin{tikzcd}
    C \ar{r}{f} \ar{d}[swap]{\varphi_t} & \P^4 \\
    C \ar{ur}[swap]{t \cdot f}
\end{tikzcd}
\end{equation*}
交换。因此，我们得到 \(\C^{\times}\) 在定义域上的作用。

我们将定义到 \(X\) 的 \(\ep\)-稳定拟映射的模空间 \(Q_{g,n}^{\ep}(X, d)\)，它类似于稳定映射的模空间。我们希望把这些模空间看作彼此虚拟双有理的，并且每一个都是带有完美障碍理论的本征 Deligne–Mumford 栈，因此我们可以定义拟映射不变量
\[ \ab< \gamma_1 \psi_1^{a_1} \cdots \gamma_n \psi_n^{a_n}>_{g,n}^{\ep} \coloneqq \int_{[Q_{g,n}^{\ep}(X,d)]^{\vir}} \prod_{i=1}^n \ev_i^*(\gamma_i) \psi_i^{a_i}. \]
其中，\(\ep \in \Q_{>0}\) 是一个正有理数。可能的 \(\ep\) 的空间具有墙与室结构，墙出现在正整数 \(n\) 的 \(\frac{1}{n}\) 处。当 \(\ep > 1\) 时我们得到 Gromov-Witten 理论（即稳定映射），而当我们向左移动时，允许的有理尾变少但基点变多。最终，在 \(\ep \to 0^+\) 极限下，不允许任何有理尾。最后，我们得到围绕 \(\frac{1}{d_0}\) 处墙的如下墙交叉公式：
\[ \ab<->_{g,n,d}^{\ep_+} - \ab<->_{g,n,d}^{\ep_-} = \sum_{k} \frac{1}{k!} \ab<-, \mu, \ldots, \mu>_{g,n+k, d-kd_0}^{\ep_+}. \]
其中，\(\mu = [zI-z]_{+, q^{d_0}, z=-\psi}\)。在特殊情形 \(g = 0\)、\(n = 1\)、\(\ep = 0^+\) 下，我们得到
\[ \sum \phi^{\alpha} \ab<\frac{\phi_{\alpha}}{z(z-\psi)}>^{\ep}_{0,1} = I_{z^{<0}}. \]
从 GW 室出发跨越所有的墙，我们得到
\[ I = 1 + \sum_{\substack{d \geq 1 \\ k \geq 0}} \frac{1}{k!} \sum_{\alpha} \phi^{\alpha} \ab< \frac{\phi_{\alpha}}{z(z-\psi)}, I_1 H - I_0 \psi, \ldots, I_1 H - I_0 \psi>_{0,1+k,d}^{\infty} q^d, \]
它通过胀子方程与除子方程产生镜面映射。

\subsection{拟映射}%
\label{sub:Quasimaps}

\begin{defn}
    到五次三维流形 \(X\) 的一个\textit{预稳定拟映射}由
    \[(C, x_1, \ldots, x_n, \mc{L}, \varphi_0, \ldots, \varphi_4),\] 
    组成，其中 \((C, x_1, \ldots, x_n)\) 是一条预稳定曲线，\(\mc{L} \in \Pic C\)，且 \(\varphi_0, \ldots, \varphi_4 \in H^0(C, \mc{L})\)，使得
    \[ \varphi^{-1}(0) \coloneqq \ab\{x \in C \mid \varphi_0(x) = \cdots = \varphi_4(x) = 0 \} \]
    有限且远离特殊点。此外，我们要求 \(\sum \varphi_i^5 = 0\)。
\end{defn}

重要的是，我们将把 \(\varphi^{-1}(0)\) 称为\textit{基点}集。在 \(\varphi^{-1}(0)\) 之外，我们有一个真正的态射 \(C \setminus \varphi^{-1}(0) \to X\)。从几何上看，假设我们在 \(\A^1 \setminus 0\) 上有一个族，使得当 \(t \to 0\) 取极限时，零轨迹 \(\varphi_i^{-1}(0)\) 最终会碰撞。我们要么在中央纤维中保留这个拟映射，要么对中央纤维作若干次爆破以得到一个真正的态射，代价是鼓出一条有理尾。因此，任何预稳定拟映射的模空间都不可能是分离的，所以我们必须考虑某种稳定性条件。

\begin{defn}
    一个预稳定拟映射 \((C, \mathbf{x}, \mc{L}, \varphi)\) 是\textit{\(\ep\)-稳定的}，如果
    \begin{enumerate}
        \item 对所有 \(x \in \varphi^{-1}(0)\)，有 \(\on{len}(x) \leq \frac{1}{\ep}\)，其中 \(\on{len}(x) = \min \ab\{ \on{ord}_x (\varphi_i)\}\)；
        \item \(\Q\)-线丛 \(\mc{L}^{\ep} \otimes \omega_C^{\log}\) 在 \(C\) 的每个不可约分支上都有正度数。
    \end{enumerate}
\end{defn}

注意到对所有不可约分支 \(C' \subseteq C\)，我们有 \(\deg \mc{L}|_{\mc{C}'} \geq 0\)。因此，发生变化的是带有 \(\leq 2\) 个特殊点的亏格 \(0\) 曲线或没有标记点的亏格 \(1\) 曲线。注意在 \((0,0)\) 与 \((1,0)\) 情形中我们有 \(C = C'\)，而在 \((0,2)\) 情形中稳定性等价于 \(\deg \mc{L} > 0\)，因此在 \((0,1)\) 情形中第二个条件等价于 
\[ \deg \mc{L}|_{C'} > \frac{1}{\ep}. \]

\begin{exm}
    在 \(\ep > 1\) 的室中，注意第一个条件意味着不存在基点，第二个条件意味着所有不稳定分支都有正度数，因此我们恢复了稳定映射。

    若我们进入 \(\frac{1}{2} < \ep < 1\) 的室，则必须收缩度数为 \(1\) 的有理尾。例如，若我们有一条以尾结尾的、由两条度数 \(1\) 有理曲线组成的链，则我们收缩最外层的尾，剩下一条带长度为 \(1\) 的基点的单度数 \(2\) 有理尾。 

    若我们进入 \(\frac{1}{3} < \ep < \frac{1}{2}\) 的室，则必须收缩所有度数 \(2\) 的有理尾，并把它们替换成长度为 \(2\) 的基点。 
\end{exm}

\begin{rmk}
    在射影空间的情形，存在一个真正的态射
    \[ Q_g^{\ep_+}(\P^N, d) \to Q_g^{\ep_-}(\P^N, d) \qquad (C_g \cup C_0, \mc{L}, \varphi) \mapsto (C_g, \mc{L}|_{C_g}(p), \varphi(p)). \]
    对一般的靶空间这做不到，但这一想法为拟映射墙交叉现象提供了直观。
\end{rmk}

我们现在讨论如何把这推广到更一般的靶空间。令
\[ X = V \sslash_{\theta} G, \]
其中 \(V\) 是仿射的，\(G\) 是约化的，\(\theta\) 是 \(G\) 的一个特征标。我们需要要求
\begin{enumerate}
    \item \(V^{\mr{ss}}(\theta) = V^s(\theta) \neq \emptyset\)；
    \item \(V\) 至多有 lci 奇点，且 \(V^s(\theta)\) 光滑。
\end{enumerate}
注意我们要求稳定轨迹具有有限的稳定化子。为方便起见，我们还要求
\begin{enumerate}
    \item \(V \sslash G\) 是射影的；
    \item \(G\) 在稳定轨迹上自由作用。
\end{enumerate}

\begin{exm}
    在五次流形的情形，\(V\) 是五次流形上的仿射锥。不稳定轨迹由原点组成，而原点是唯一的奇点。
\end{exm}

\begin{exm}
    另一个例子是，我们可以考虑带有一般稳定性条件的 \(\C^N \sslash (\C^{\times})^r\)，由此得到一个环簇。
\end{exm}

\begin{exm}
    我们可以考虑 Grassmannian
    \[ G(2,4) = M_{2\times 4} \sslash GL_2, \]
    其中不稳定轨迹由秩严格小于 \(2\) 的矩阵组成。
\end{exm}

\begin{rmk}
    当 \(G\) 是交换群时，对一般的稳定性条件，稳定轨迹上的条件总是满足的，但对非交换的 \(G\) 这并不总能做到。
\end{rmk}

\begin{rmk}
    对一般的光滑射影簇，原点未必是 lci 的，因此如何定义拟映射且仍能定义虚拟闭链并不清楚。
\end{rmk}

\begin{defn}
    到 \(V \sslash G\) 的一个\textit{预稳定拟映射}是一个映射 \(u \colon C \to [V/G]\)，使得基轨迹
    \[ u^{-1}([V^{\mr{un}}(\theta) / G]) \]
    有限且远离特殊点。按定义，\(u\) 由一个主 \(G\)-丛 \(P \to C\) 连同一个等变映射 \(P \to V\) 组成，这等价于相伴丛 \(P \times_G V\) 的一个截面。 
\end{defn}

\begin{exm}
    对五次流形而言，这恰好恢复了一个线丛 \(\mc{L} \in \Pic C\)（等价于一个 \(\C^{\times}\)-挠子）连同 \(\on{Tot}(\mc{L}^{\oplus 5})\) 的一个满足定义方程的截面。
\end{exm}

为了定义度数，我们可以考虑 \([V/G]\) 上的线丛 \(\mc{L}_{\theta}\)。于是我们可以定义 \(C' \subset C\) 在 \(u\) 下的度数为
\[ \deg (u^* \mc{L}_{\theta}|_{C'}). \]

\begin{exm}
    在 \(G(2,4)\) 的例子中，注意一个主 \(\on{GL}_2\)-丛等价于 \(C\) 上的一个秩 \(2\) 向量丛 \(\mc{E}\)。因此，\(P \times_G V\) 等价于 \(\mc{E}^{\oplus 4}\) 的数据，从而到 Grassmannian 的拟映射等价于稳定商。在这一情形中，我们要找的是在一般点处满射的
    \[ \mc{O}^{\oplus 4} \to \mc{E}, \]
    并且因为 \(\theta = \det\)，我们看到 \(u^* \mc{L}_{\theta} = \det \mc{E}\)。
\end{exm}

\begin{exm}
    考虑 \(V = \C^{N+1}\)，并注意 \(\P^N = V \sslash \C^{\times}\)。若考虑 \(\theta_1 \colon z \mapsto z\)，则有 
    \[ V^{\on{un}} = V((x_i)). \]
    然而，若考虑 \(\theta \colon z \mapsto z^2\)，则有
    \[ V^{\on{un}} = V((x_i x_j)). \]
\end{exm}

这里我们看到，概形结构对特征标的选择很敏感，因此不能简单地朴素定义基点的长度。注意到
\[ V^{\on{un}} = V\ab(\ab< \bigoplus_{k \geq 1} R_k>), \qquad R_k = H^0([V/G], L_{\theta}^{\otimes k}). \]
特别地，我们必须定义
\[ \on{len}(x) \coloneqq \min_{\substack{f \in R_k \\ k \geq 1}} \frac{\on{ord}_x u^*(f)}{k}. \]

\begin{exm}
    令 \(X = \mr{pt} = \C \sslash_{\theta} \C^{\times}\)。此时一个预稳定拟映射就是
    \[ (C, \mc{L}, \varphi \in H^0(\mc{L})), \]
    其中 \(\varphi^{-1}(0)\) 有限且远离特殊点。此外，若 \(\deg \mc{L} = 1\)，我们恰好恢复了 \(C^{\mr{sm}}\) 中一个新标记点的数据。
\end{exm}

\begin{exm}
    现在考虑 \(X = \mr{pt} = \C^n \sslash (\C^{\times})^n\)。同样设度数为 \((1, \ldots, 1)\)。现在我们令
    \[ \theta = (a_1, \ldots, a_n). \]
    因此注意到
    \[ \ab<R_1> = \ab<(x_1^{a_1} \cdots x_n^{a_n})>. \]
    因为它们都处在度数 \(1\) 中，我们看到对光滑点 \(x \in C^{\mr{sm}}\) 有
    \[ \on{len}_x = \sum \on{ord}_x (\varphi_i^{a_i}) = \sum a_i \on{ord}_x \varphi_i. \]
    若记 \(\ab\{p_i\} = \varphi_i^{-1}(0)\)，则条件 \(\ep \cdot \on{len}(x) \leq 1\) 等价于要求 
    \[ \sum_{x = p_i} \ep a_i \leq 1. \]
    度数条件现在说的是
    \[ \ep \sum_i a_i + \deg \omega_C > 0. \]
    等价地，这说的是
    \[ \omega_C(\ep a_1 p_1 + \cdots + \ep a_n p_n) > 0, \]
    于是我们恢复了一个权为 \((\ep a_1, \ldots, \ep a_n)\) 的加权标记曲线。
\end{exm}

\subsection{主空间}%
\label{sub:Master space}

证明使用虚拟局部化~\cite{GP1999}；Hassett 的加权模空间~\cite{Hassett2003} 为这一构造提供了模型之一。

我们的方法是在一个插值两个模空间的主空间上作虚拟局部化。考虑 \(\P^n\) 上由
\[ t \cdot [x_0 : \cdots : x_n] = [t x_0, x_1, \ldots, x_n]. \]
给出的 \(\C^{\times}\)-作用。不动点轨迹由一个 \(x_0 = 0\) 处的 \(\P^{n-1}\) 和一个点 \([1, 0, \ldots, 0]\) 组成。于是我们可以通过把计算转移到另一个不动点轨迹来算出
\[ \int_{\P^{n-1}} H^{n-1} = 1 \]
特别地，我们把它提升为 \(H \in H_{\C^{\times}}^* (\P^n) = \Q[H,z] / (H-z)H^n\) 并计算
\begin{align*}
    \int_{\P^n} H^{n-1} &= \int_{\P^{n-1}} \frac{H^{n-1}}{-z+H} + \int_{\mr{pt}} \frac{z^{n-1}}{z^n}.
\end{align*}
取留数，我们看到
\begin{align*}
    0 &= \on{Res}_{z = 0}\int_{\P^{n-1}} \frac{H^{n-1}}{-z+H} + 1 \\
    &= 1 - \int_{\P^{n-1}}H^{n-1},
\end{align*}
这恢复了所需的积分。我们这里用到的关键事实是射影空间是本征的。

\begin{exm}[第二个玩具例子]
    我们将考虑第二个玩具例子。回忆 Hassett 的加权标记曲线模空间
    \[ (C, x_1, \ldots, x_n) \]
    其中 \(C\) 是结点曲线，且 \(x_i \in C^{\mr{sm}}\) 可以碰撞。碰撞由如下方式调控：每个标记点被赋予一个权 \(\delta_i \in \Q \cap (0,1]\)，并要求满足以下条件：
    \begin{enumerate}
        \item 对所有 \(x \in C^{\mr{sm}}\)，有
            \[ \sum_{x = x_i} \delta_i \leq 1; \]
        \item 我们有稳定性条件 \(\omega_C(\delta_1 x_1 + \cdots + \delta_n x_n) > 0\)。例如，若我们有一条有理尾，则 \(\sum_{x_i \in C'} \delta_i > 1\)。
    \end{enumerate}
\end{exm}

\begin{exm}[又一个玩具例子]
    我们的目标是再次计算
    \[ \int_{\bar{M}_{0,4}} \psi_1 = 1. \]
    我们的第一个目标是改变稳定性条件。若令
    \[ \delta_+ = \ab(\frac{2}{3}, \frac{2}{3}, \frac{2}{3}, \frac{2}{3}), \]
    我们看到稳定性条件没有改变。然而，若令
    \[ \delta_- = \ab(\frac{2}{3}, \frac{2}{3}, \frac{2}{3}, \frac{1}{3}), \]
    模空间没有改变，但万有曲线实际上变成了 \(\P^1 \times \P^1\)，其中截面 \(x_4\) 是对角线。这告诉我们
    \[ \int_{\bar{M}_{0,4}^{\delta_-}} \psi_1 = 0, \]
    并且我们将利用它在 \(\delta_+\) 室中作计算。我们可以尝试构造一个主空间：既允许带有 \(x_i, x_4\) 的有理尾，又允许 \(x_1 = x_4\)；但若这样做，我们最终会得到一个半稳定点，它带有一条满足 \(x_1 = x_4\) 的有理尾。我们希望通过对一个主空间取商来得到它，因此我们将给模空间添加一个切向量的数据。于是，定义 \(\delta_{\pm}\) 主空间由形如
    \[ \xi = (C, x_1, \ldots, x_4, v \in T_{x_4} C \cup \infty) \]
    的对象组成，其中 \(C\) 是结点曲线，\(x_i\) 是光滑点，\(x_1, x_2, x_3\) 互不相交，
    \[ \sum_{x_i \in C'} \frac{2}{3} \delta_i > 1 \]
    对所有有理尾 \(C'\) 成立，若 \(v = 0\)，则 \(x_4\) 与其余 \(x_i\) 不相交，而当 \(v = \infty\) 时，不存在满足 \(x_4 \neq x_i\)（\(i = 1,2,3\)）的有理尾。 

    我们现在通过缩放 \(v\) 给这个主空间一个 \(\C^{\times}\) 作用。有若干不动点轨迹：
    \begin{enumerate}
        \item 当 \(v = 0\) 时，我们得到 \(\bar{M}_{0,4}^{\delta_+}\)；
        \item 当 \(v = \infty\) 时，我们得到 \(\bar{M}_{0,4}^{\delta_-}\)；
        \item 还有一些不动点轨迹，其中 \(v \neq 0,\infty\) 且我们有一条满足 \(x_i = x_4\) 的有理尾。
    \end{enumerate}

    若考虑一个一般点（四个标记点互不相交），它将沿 \(\C^{\times}\) 作用直接行到自身。然而，若我们从一条包含 \(x_1\) 与 \(x_4\) 的有理尾出发，则我们先走到一条 \(v\) 一般但 \(x_1 = x_4\) 的有理尾，而若把结点光滑化，我们得到一条 \(x_1 = x_4\) 的单曲线，但我们允许 \(v \to \infty\)，这是另一个不动点。现在我们可以积分
    \begin{align*}
        \int_{\ms{Master}} \psi_1 = \int_{\bar{M}_{0,4}^{\delta_+}} \frac{\psi_1}{-z-?} + \int_{\bar{M}_{0,4}^{\delta_-}} \frac{\psi_1}{z+?} + \on{Cont}(12|34) + \on{Cont}(13|24) + \on{Cont}(14|23).
    \end{align*}
    在 \(z = 0\) 处取留数，我们得到
    \begin{align*}
        0 &= - \int_{\bar{M}_{0,4}^{\delta_+}} \psi_1 + 0 + 0 + 0 + \int_{\mr{pt}} \frac{(-z)}{(\pm z)(\pm z)},
    \end{align*}
    其中最后一项分母中的 \(\pm z\) 是结点光滑化参数，因此我们恢复了所需的积分。
\end{exm}

有了所有这些玩具例子，我们终于准备好讨论拟映射主空间。我们在 \(\frac{1}{d_0}\) 墙处的第一步是同时允许度数 \(d_0\) 的有理尾和长度 \(d_0\) 的基点。然而，我们可以考虑带有一条有理尾且拟映射形如
\[ [a_0 u^{d_0}, \ldots, a_n u^{d_0}], \]
的曲线，它们有无穷多个自同构。然而，我们不能在基点处附加切向量（因为形变会消除它们），并且还存在这样的情形：有多于一条度数 \(d_0\) 的尾带有一个长度为 \(d_0\) 的基点（此时自同构群包含 \(\C^{\times}\)）。

\begin{rmk}
    有一个关键的观察。假设我们有两条度数 \(d_0\) 的有理尾，并把两个结点记为 \(p_1\) 与 \(p_2\)。现在定义光滑化参数空间
    \[ \C \cong \Theta_i = T_{p_i} E_i \otimes T_{p_i} C_g, \]
    其中 \(E_i\) 是有理尾，\(C_g\) 是曲线的其余部分。我们将考虑按度数标记的加权半稳定曲线的模空间
    \[ \mf{M}_{g,n,d}^{\mr{wt}} \xrightarrow{\text{\'et}} \mf{M}_{g,n} \]
    注意这个态射是高度非分离的。我们将考虑开子栈 \(\mf{M}_{g,n,d}^{\mr{wt}, \ep_0}\)，其中的曲线没有度数严格小于 \(d_0\) 的尾。于是我们可以考虑 Cartier 除子 \(\mc{Z}\)，即至少存在一条度数 \(d_0\) 的尾的轨迹。特别地，我们看到
    \[ \Theta_i = N_{\mc{Z}_i} |_{\xi} \]
    其中 \(\mc{Z}_1\) 与 \(\mc{Z}_2\) 是所关心点处的两个分支。我们还看到 \(\Theta_1 \otimes \Theta_2\ = \mc{O}(\mc{Z})|_{\xi}\)，因此我们可以添加
    \[ v \in \ab(\bigotimes_i \Theta_i)^{\vee} \cup \ab\{\infty\} \]
    的数据以及 
    \[ w \in \P(\Theta_1 \oplus\Theta_2). \]
    的数据。
\end{rmk}

实际的构造如下进行。考虑 \(\mc{Z} \subset \mf{M} = \mf{M}_{g,d}^{\mr{wt}, \ep_0}\)。定义 \(\tilde{\mf{M}} \to \cdots \to \mf{M}\) 为沿 \(\mc{Z}\) 的层片的迭代爆破。于是主空间（带校准尾的拟映射）在概形 \(S\) 上的一个对象由如下数据给出：一个拟映射
\begin{equation*}
\begin{tikzcd}
    C \ar{d} \ar[dashrightarrow]{r}{f} & X \\
    S
\end{tikzcd}
\end{equation*}
它诱导一个态射 \(p \colon S \to \mf{M}_{g,d}^{\mr{wt}, \ep_0}\)，一个提升的选取
\begin{equation*}
\begin{tikzcd}
    & \tilde{\mf{M}} \ar{d}{\tau} \\
    S \ar{r}{p} \ar[dashrightarrow]{ur}{w} & \mf{M}_{g,d}^{\mr{wt}, \ep_0},
\end{tikzcd}
\end{equation*}
一个线丛 \(\mc{N} \in \Pic S\)，以及 \(v = (v_1, v_2)\)，其中 \(v_1 \in H^0(w^* \mc{O}(-\mc{Z}))\) 且 \(v_2 \in H^0(\mc{N})\)，且 \((v_1, v_2)\) 处处非零。我们还不允许长度严格大于 \(d_0\) 的基点，并要求如下的稳定性：
\begin{enumerate}
    \item 当 \(v = 0\) 时，我们得到 \(\ep_+\)-稳定性；
    \item 当 \(v = \infty\) 时，我们得到 \(\ep_-\)-稳定性；
    \item 任何常值尾（度数 \(d_0\) 且带有长度 \(d_0\) 的基点的尾）都是纠缠尾。这里，我们考虑最后一次爆破，使得 \(w\) 映到一个例外轨迹 \(\bigcap_{i \in I} \tilde{\mc{Z}}_i\) 中，此时这些 \(E_i\)（\(i \in I\)）就是纠缠尾；
    \item 若存在度数 \(d_0\) 的尾，则所有长度 \(d_0\) 的基点都位于度数 \(d_0\) 的尾上。
\end{enumerate}






